\chapter*{Abstract}
\addcontentsline{toc}{chapter}{Abstract}
\markboth{Abstract}{Abstract}

\noindent

Virtual orbital cost is a fundamental bottleneck in post-Hartree–Fock quantum chemistry and near-term quantum computing. Computing correlation energy via MP2, coupled cluster, or variational quantum algorithms requires expensive summations over all virtual orbitals, severely limiting the molecular system size and circuit depth on NISQ devices. This thesis implements the Optimized Virtual Orbital Space (OVOS) method of Adamowicz \& Bartlett (1987). The approach minimises a second-order Hylleraas functional via unrestricted Hartree–Fock reference states and Newton–Raphson orbital rotations. Validation across five representative molecules  (H$_2$O, CO, HF, NH$_3$, Li$_2$)  in both 6--31G and cc--pVDZ basis sets demonstrates that ~90\% of full-space MP2 correlation energy is recovered using 50–68\% of virtual orbitals, except Li$_2$. Li$_2$ shows greater consistency with the original Adamowicz \& Bartlett reference calculations at 28\% of virtual orbitals. However, convergence becomes challenging as the active space approaches full dimension. OVOS optimised orbitals provide superior trial states for variational quantum eigensolvers compared to canonical Hartree–Fock, reducing circuit depth while maintaining energy accuracy. These findings demonstrate practical value for NISQ-era quantum chemistry.

\bigskip

\textbf{Keywords:} Optimized virtual orbital space, OVOS, MP2, orbital optimisation,
virtual orbital reduction, quantum computing, PySCF, Newton-Raphson.

